\section{Completion of the construction and the Ramsey bounds}
\label{end:section}

The preceding compression step replaces an old description by a new
one while retaining a fixed positive fraction of the selected tuple.
We now iterate that replacement a bounded number of times. The final
description has too little entropy to encode a selected tuple from the
original stream. We then pass from the resulting prime-indexed graphs
to every sufficiently large integer $t$ and prove the matching upper
estimate.

\subsection{A bounded number of compression stages}
\label{end:construction}

\begin{proof}[Proof of Theorem~\ref{intro:prime-construction}]
Fix $d\ge5$ and $0<\eta<1/10$, with all the parameters of
Equation~\eqref{base:stream-parameters}. Suppose, for a contradiction,
that every stream has a consistent $k$-tuple. Select one by a fixed
rule, and condition on the simultaneous rectangle event of
Lemma~\ref{base:rectangles}. For large $q$ this event has probability
at least $1/2$, so the new stream law is dominated by twice the iid
law. The selected tuple has deterministic length $k$.

Initially the context is empty. Its entropy budget is $\Lambda_0=0$,
and the full flag domain has size
\[
 Q_dQ_{d-1}\le4q^d\exp((d-1)\sigma).
\]
Thus we may take $\Delta_0=(d-1)\sigma$. The stage parameter defined
in Equation~\eqref{mark:parameters} is initially
\[
 D_0=\sigma^\beta\bigl(1+(d-1)\sigma^{1-\eta}\bigr)
     =O(\sigma^{1-\eta+\beta}),
\]
which lies in the admissible interval
$[\sigma^\beta,\sigma^{1-\eta/2}]$ for sufficiently large $q$.

Consider any admissible stage with budgets $\Lambda,\Delta$ and
parameter $D$. By Lemma~\ref{comp:step}, its replacement context has
budgets satisfying
\begin{equation}
 \Lambda'\le CqP\log(w+2)(\sigma+wP)+Cw\sigma,
 \qquad \Delta'\le CK_*,
 \qquad 1\le w\le C\sigma^{1+\eta/2}/D.
 \label{end:one-step-budget}
\end{equation}
The high-rank case is included with $w=1$. The constants may depend on
$d,\eta$, a fixed lower bound for $\ell/k$, and fixed upper bounds for
the stream domination factor and the domain prefactor. They are uniform
over all input laws, contexts, conditional histories, and class choices
satisfying those numerical bounds, as in Lemma~\ref{comp:step}.
Using $k\asymp q\sigma^{1+\eta}$ and
$P\le4D\sigma^{9\beta}$, the three terms to estimate satisfy
\[
 \frac{qP\sigma}{k}\le CD\sigma^{9\beta-\eta},\qquad
 \frac{qP^2w}{k}\le CD\sigma^{18\beta-\eta/2},\qquad
 \frac{w\sigma}{k}\le\frac{Cq^{-1}\sigma^{1-\eta/2}}{D}.
\]
Also $\log(w+2)=O(\log(\sigma+2))$. We obtain
\begin{align}
 \frac{\Lambda'}{k}
 &\le CD\log(\sigma+2)
       \bigl(\sigma^{9\beta-\eta}
                  +\sigma^{18\beta-\eta/2}\bigr)
       +\frac{Cw\sigma}{k}\notag\\
 &\le D\sigma^{-\eta/3}.
 \label{end:entropy-contraction}
\end{align}
Indeed, $w\sigma/k\le Cq^{-1}\sigma^{1-\eta/2}/D$ is exponentially
small in $\sigma$, uniformly for admissible $D$. The other two terms
have exponents strictly below $-\eta/3$, because
$\beta=\eta/10^7$; their fixed constants and logarithmic factor are
absorbed for large $q$.

Recompute the stage parameter from the new budgets, without retaining
the old context. Equations~\eqref{mark:parameters} and
\eqref{end:entropy-contraction} give
\begin{align}
 D_{\rm new}
 &=\sigma^\beta\bigl(1+\Lambda'/k+\Delta'\sigma^{-\eta}\bigr)
   \notag\\
 &\le\sigma^\beta
     +D\sigma^{\beta-\eta/3}+CD\sigma^{7\beta-\eta}
   \notag\\
 &\le\sigma^{2\beta}+D\sigma^{-\eta/4}.
 \label{end:parameter-contraction}
\end{align}
The last inequality follows from
$\beta-\eta/3<-\eta/4$ and $7\beta-\eta<-\eta/4$.
This recurrence preserves admissibility: its lower bound
$D_{\rm new}\ge\sigma^\beta$ follows from the definition, while
\[
 \sigma^{2\beta}+\sigma^{1-3\eta/4}
 \le\sigma^{1-\eta/2}
\]
for large $q$.

Choose the number of stages in advance, for example
$T=\lceil8/\eta\rceil$. Iterating
Equation~\eqref{end:parameter-contraction}, and using $D_0\le\sigma$
for large $q$, gives
\[
 D_T\le\sigma^{1-T\eta/4}
       +\frac{\sigma^{2\beta}}{1-\sigma^{-\eta/4}}
       \le2\sigma^{2\beta}.
\]
One further compression stage produces budgets $\Lambda_*,\Delta_*$
with
\begin{equation}
 \frac{\Lambda_*}{k}
     \le2\sigma^{2\beta-\eta/3}=o(1),
 \qquad
 \frac{\Delta_*q\sigma}{k}
     \le C\sigma^{8\beta-\eta}=o(1).
 \label{end:final-budgets}
\end{equation}

We justify explicitly the probabilistic invariants throughout this
iteration. At each stage Lemma~\ref{comp:step} first averages over
fresh public tables and fixes a table family with an extraction event
of probability bounded below. It then conditions on that event and
retains a prescribed deterministic number of slots in one of the two
allowed orientations. The new decoder uses only its new context and
its fixed tables. In particular, its entropy budget is $\Lambda'$,
not the sum of earlier context entropies. Every retained tuple is
consistent, remains in the original stream in an allowed orientation,
and still satisfies the simultaneous rectangle bound.

The constants and thresholds can be fixed without knowing the actual
conditional laws. Start with the numerical bounds $\ell/k=1$, domain
prefactor $4$, and domination factor $2$. The uniform statement of
Lemma~\ref{comp:step} supplies a positive lower bound on the next
retained proportion and on the conditioning probability, and an upper
bound on the next domain prefactor. Use these to prescribe the next
stage's numerical bounds, including the old domination bound divided
by the success-probability lower bound. Recursively do this for the
already fixed $T+1$ stages, taking maxima over the finitely many class
types if necessary. All resulting constants are independent of $q$;
choose $q$ above their finitely many uniform thresholds and the
arithmetic thresholds used above. No enumeration of conditional
histories, or history-dependent limiting argument, is involved.

Consequently the final length is deterministic and satisfies
$\ell_*=\Theta(k)$, and the final stream law still has a constant
domination factor. Lemma~\ref{base:selection-entropy} therefore
applies to the final tuple $F_*$. Only the marking scans are now
performed: no further class selection, tuple extraction, or
conditioning is used. The marking message specifies the expensive
flag values and a cheap domain for each remaining slot. Conditional
on that message, the entropy of the remaining values is bounded by
the logarithm of the product of those domains. Thus
Equation~\eqref{mark:entropy-upper} bounds the entropy of the entire
tuple under its present law.

Therefore, Equation~\eqref{mark:entropy-upper} from the final
marking scan and Equation~\eqref{end:final-budgets} give
\[
 H(F_*)\le d\sigma\ell_*+\Lambda_*+O(k)+O(\Delta_*q\sigma)
          =d\sigma\ell_*+O(k).
\]
Lemma~\ref{base:selection-entropy} gives instead
\[
 H(F_*)\ge d\sigma\ell_*+\eta\ell_*\log\sigma-O(k).
\]
These inequalities contradict each other, since
$\ell_*\ge c k$ for a fixed $c>0$ and $\log\sigma\to\infty$.
Consequently some stream has no consistent $k$-tuple. Its flag graph
has independence number less than $k$ and is $K_{d+1}$-free by
Lemma~\ref{base:flag-graph}, as required.
\end{proof}

\subsection{From prime orders to every integer}

Only a fixed-ratio interval between consecutive available prime scales
is needed. We include an elementary proof to avoid imposing a
prime-subsequence restriction on the Ramsey bound.

\begin{lemma}[A prime in a fixed-ratio interval]
\label{end:prime-interval}
There is a constant $c_0>0$ such that, for every sufficiently large real
$x$, a prime $q$ satisfies $c_0x\le q\le x$.
\end{lemma}

\begin{proof}
Define
\[
 \vartheta(x)=\sum_{p\le x}\log p,
 \qquad
 \psi(x)=\sum_{\substack{p^h\le x\\h\ge1}}\log p,
\]
where $p$ ranges over primes. Every prime in $(m,2m]$ divides
$\binom{2m}{m}$, so
\[
 \vartheta(2m)-\vartheta(m)
 \le\log\binom{2m}{m}\le2m\log2.
\]
Summing at powers of two and then using monotonicity proves
$\vartheta(x)\le Cx$.

The exponent of $p$ in the central binomial coefficient is
\[
 \sum_{h\ge1}
 \left(\left\lfloor\frac{2m}{p^h}\right\rfloor
        -2\left\lfloor\frac{m}{p^h}\right\rfloor\right).
\]
Each summand is zero or one, and therefore
\[
 \psi(2m)\ge\log\binom{2m}{m}
 \ge2m\log2-\log(2m+1).
\]
The last inequality follows because the largest of the $2m+1$
binomial coefficients is at least their average $4^m/(2m+1)$.
For proper prime powers the rough estimate
\[
 0\le\psi(x)-\vartheta(x)
   \le C\sqrt{x}\log^2(x+2)=o(x)
\]
suffices: there are $O(\log x)$ exponents $h\ge2$, at most
$\sqrt{x}$ possible bases for each, and each weight is at most
$\log x$. It follows first for even integers, and then by monotonicity
for all large real $x$, that $\vartheta(x)\ge cx$.
Choose $c_0<c/(2C)$. For sufficiently large $x$,
$\vartheta(x)>\vartheta(c_0x)$, proving that the stated interval
contains a prime.
\end{proof}

Fix $s\ge6$, put $d=s-1$, and choose $0<\eta<1/10$.
For a sufficiently large integer $t$, apply
Lemma~\ref{end:prime-interval} with
\[
 x=\frac{t}{(\log t)^{1+\eta}}.
\]
The resulting prime tends to infinity, and $\log q\sim\log t$.
Moreover,
\[
 q(\log q)^{1+\eta}
 \le t\left(\frac{\log q}{\log t}\right)^{1+\eta}<t.
\]
The graph from Theorem~\ref{intro:prime-construction} therefore has no
independent set of size $t$. Its order gives
\begin{equation}
 r(s,t)>\lfloor q^d\log q\rfloor
 \ge c_{s,\eta}\frac{t^d}{(\log t)^{d-1+d\eta}}.
 \label{end:lower-slack}
\end{equation}
Here the floor is absorbed into a fixed positive constant for large
$t$. Given any $\varepsilon>0$, choose
$\eta=\min\{1/20,\varepsilon/(2d)\}$. Then
$\varepsilon-d\eta>0$, so
$(\log t)^{\varepsilon-d\eta}$ eventually absorbs
$1/c_{s,\eta}$. Equation~\eqref{end:lower-slack} yields the required
coefficient-one lower bound for every sufficiently large integer $t$.

\subsection{The matching upper bound}

The upper estimate of Ajtai, Koml\'os, and Szemer\'edi \cite{AKS1980}
can be obtained from a triangle-free independence bound followed by
sampling and deletion. Shearer proved a lower bound asymptotic to
$n\log d/d$ as $d\to\infty$ for triangle-free graphs of average degree $d$
\cite{Shearer1983}. The elementary proof below uses a uniformly chosen
independent set, following Alon's method
\cite[proof of Proposition~2.1]{Alon1996}; sharper
occupancy estimates appear in Davies, Jenssen, Perkins, and Roberts
\cite[Theorems~1 and~3]{DJPR2018}. We retain a convenient constant and give the details
separately from the lower-bound construction.

\begin{lemma}[Independent sets in a triangle-free graph]
\label{end:triangle-free}
If a triangle-free graph $H$ has $n>0$ vertices and maximum degree at
most a real number $D_0\ge1$, then
\[
 \alpha(H)\ge\frac{n\log(D_0+1)}{8(D_0+1)}.
\]
\end{lemma}

\begin{proof}
Choose an independent set $I$ uniformly from all independent sets of
$H$, including the empty set, and write $\mu=\E|I|/n$.
At a vertex $x$, expose $I$ outside the closed neighborhood of $x$.
Let $Z_x$ be the number of neighbors that have no neighbor in this
exposed part of $I$. These available neighbors are mutually
nonadjacent because $H$ is triangle-free. The remaining independent-set
choices are therefore $\{x\}$ or an arbitrary subset of the $Z_x$
available neighbors, all with equal conditional probability. Hence
\[
 \Prob(x\in I\mid\text{exposure})=\frac1{1+2^{Z_x}},
 \qquad
 \E(|I\cap N(x)|\mid\text{exposure})
   =\frac{Z_x2^{Z_x-1}}{1+2^{Z_x}}\ge\frac{Z_x}{4}.
\]
The last inequality includes $Z_x=0$.
Averaging also over a uniform vertex $x$, the degree bound gives
\[
 \frac{\E Z_x}{4}
 \le\frac1n\E\sum_{v\in I}\deg(v)\le D_0\mu.
\]
Since $(1+2^z)^{-1}\ge2^{-z}/2$ for $z\ge0$, Jensen's inequality
now gives
\begin{equation}
 \mu\ge\tfrac12\E 2^{-Z_x}
       \ge\tfrac12 2^{-\E Z_x}
       \ge\tfrac12 2^{-4D_0\mu}.
 \label{end:triangle-jensen}
\end{equation}
If $\mu<\log(D_0+1)/(8(D_0+1))$, then
Equation~\eqref{end:triangle-jensen} implies
$\mu>1/(2\sqrt{D_0+1})$, using $\log2<1$.
This is impossible: $\log u\le2\sqrt u$ for $u\ge1$ implies
$\log u/(8u)<1/(2\sqrt u)$. Thus $\mu$ has the claimed lower
bound, and $\alpha(H)\ge\E|I|=n\mu$ completes the proof.
\end{proof}

\begin{proposition}[Upper Ramsey estimate]
\label{end:upper-bound}
For every fixed integer $s\ge2$ there is a constant $C_s$ such that
for all sufficiently large integers $t$,
\[
 r(s,t)\le C_s\frac{t^{s-1}}{(\log t)^{s-2}}.
\]
\end{proposition}

\begin{proof}
Induct on $s$. The case $s=2$ follows from $r(2,t)=t$.
If a triangle-free graph has no independent set of size $t$, then all
its degrees are below $t$, because every neighborhood is independent.
Lemma~\ref{end:triangle-free}, with $D_0=t$, bounds its order by
$8t(t+1)/\log(t+1)$. This proves the case $s=3$.

Now let $s\ge4$ and suppose the estimate holds for smaller clique
parameters. Let $G$ have $n>0$ vertices, maximum degree $h$, no $K_s$,
and no independent set of size $t$. Every neighborhood avoids
$K_{s-1}$ and an independent $t$-set. Likewise the common neighborhood
of the two endpoints of an edge avoids $K_{s-2}$ and an independent
$t$-set. Therefore
\[
 h<r(s-1,t),\qquad
 |N(u)\cap N(v)|\le m:=r(s-2,t)\quad(uv\in E(G)).
\]
Summing these common-neighborhood sizes over edges counts each
triangle three times, so $G$ has at most $nhm/6$ triangles.

If $h\le t^{s-5/2}$, the greedy bound
$\alpha(G)\ge n/(h+1)$ gives
\[
 n<t(h+1)=O(t^{s-3/2})
   =o\!\left(\frac{t^{s-1}}{(\log t)^{s-2}}\right).
\]
Assume henceforth that $h>t^{s-5/2}$, and set
$\lambda=(hm)^{-1/2}\le1$. By the induction bound on $m$,
\begin{equation}
 \lambda h=\sqrt{h/m}
 \ge c_s t^{1/4}(\log t)^{(s-4)/2}\ge t^{1/5}
 \label{end:sample-degree}
\end{equation}
for sufficiently large $t$.

Retain each vertex independently with probability $\lambda$.
Delete all retained vertices whose degree in the sampled graph exceeds
$12\lambda h$, and then delete one vertex from each triangle still
present until none remains. If $E_0,T_0$ are the initial sampled edge
and triangle counts, respectively, the first deletion removes at most
$2E_0/(12\lambda h)$ vertices, and the second at most $T_0$.
Their expected numbers are at most
\[
 \frac{\lambda^2nh}{12\lambda h}=\frac{n\lambda}{12},
 \qquad
 \frac{\lambda^3nhm}{6}=\frac{n\lambda}{6}.
\]
Some realization therefore leaves a triangle-free induced subgraph
$H$ with
\[
 |V(H)|\ge\tfrac34n\lambda,
 \qquad \Delta(H)\le12\lambda h.
\]
Deleting further vertices cannot increase degrees, so both conclusions
hold simultaneously.
Lemma~\ref{end:triangle-free} and
Equation~\eqref{end:sample-degree} imply
\[
 t>\alpha(H)
 \ge\frac{(3/4)n\lambda\log(12\lambda h+1)}{8(12\lambda h+1)}
 \ge c\frac{n\log t}{h}.
\]
Finally substitute $h<r(s-1,t)$ and its induction bound to obtain
$n\le C_s t^{s-1}/(\log t)^{s-2}$. Increasing $C_s$ absorbs the
additive one in passing from the order of an avoiding graph to its
Ramsey number.
\end{proof}

\subsection{The logarithmic exponent}
\label{end:main-proof}

\begin{proof}[Proof of Theorem~\ref{intro:main}]
The lower bound follows from Equation~\eqref{end:lower-slack} and the
choice of $\eta$ made there. Proposition~\ref{end:upper-bound} gives
the asserted upper estimate. For every fixed $\varepsilon>0$, its
constant is eventually bounded by $(\log t)^\varepsilon$; hence
\[
 \frac{t^{s-1}}{(\log t)^{s-2+\varepsilon}}
 \le r(s,t)\le
 \frac{t^{s-1}}{(\log t)^{s-2-\varepsilon}}
\]
for all sufficiently large integer $t$. Taking logarithms and dividing
by $\log\log t>0$ bounds the expression in the theorem between
$s-2-\varepsilon$ and $s-2+\varepsilon$. Since $\varepsilon$ is
arbitrary, its limit is $s-2$.
\end{proof}
